\section{Evaluation}
\label{sec:evaluation}

\subsection{Experimental Setup}

We run VRASection and GeoSection on seven states with historically significant
VRA litigation or demographic composition at the 2020 Census.
All runs use the 2020 TIGER census-tract adjacency graphs with TIGER boundary
lengths as edge weights.
Minority VAP is sourced from the 2020 Census P4 table (Hispanic or Latino,
and Not Hispanic or Latino -- Black or African American alone).
State congressional allocations are the 2020 apportionment (in effect for
the 2022 elections).
VRASection parameters: $N = 50$ seeds per ratio, $\delta = 0.005$ balance
tolerance, $w_\text{vra} = 0.40$ alignment weight.

\begin{table}[t]
\centering
\caption{VRA evaluation states: demographics and litigation history.}
\label{tab:states}
\begin{tabular}{llcll}
\toprule
\textbf{State} & \textbf{$k$} & \textbf{Minority VAP (\%)} & \textbf{Group} & \textbf{Key cases} \\
\midrule
Alabama (AL) & 7  & 27.5 & Black      & \emph{Allen v.\ Milligan} (2023) \\
Mississippi (MS) & 4 & 37.6 & Black   & \emph{NAACP v.\ Fordice} (1992) \\
Georgia (GA) & 14 & 33.0 & Black+Hisp & \emph{Georgia v.\ Ashcroft} (2003) \\
North Carolina (NC) & 14 & 23.0 & Black & \emph{Thornburg v.\ Gingles} (1986) \\
Louisiana (LA) & 6 & 33.2 & Black     & \emph{Callais v.\ Landry} (2025) \\
South Carolina (SC) & 7 & 27.5 & Black & \emph{Alexander v.\ SC NAACP} (2023) \\
\bottomrule
\end{tabular}
\end{table}

\subsection{Multi-State Comparison}

Table~\ref{tab:comparison} reports the first-level bisection ratio selected by
GeoSection (B.8) and VRASection (this paper) for each state, along with
MetisVra (B.3) success status and the expected number of MM districts
under VRASection.

\begin{table}[t]
\centering
\caption{GeoSection vs.\ VRASection first-level bisection ratios across VRA states.
  MetisVra success column is from B.3 \citep{b3paper}.
  $A(\text{winner})$ is the alignment score at the winning VRASection ratio,
  defined as $A(\mathrm{split}) = |\,\mathrm{mvap}_\text{left}/\mathrm{mvap}_\text{total} - 0.5\,| \times 2$.
  ``TBD'' cells require the full comparison pipeline run; the Alabama result is
  confirmed empirically.}
\label{tab:comparison}
\begin{tabular}{lccccc}
\toprule
\textbf{State} & \textbf{GeoSection ratio} & \textbf{VRASection ratio}
  & \textbf{$A(\text{winner})$}
  & \textbf{MetisVra success?} & \textbf{VRASection MM districts} \\
\midrule
AL & 3:4 & \textbf{1:6} & 0.27 & No (0/28 configs) & Black Belt isolated \\
MS & 1:3 & \textbf{1:3} & 0.12 & Yes (23/28) & Same (minority dispersed) \\
GA & 1:13 & \textbf{1:13} & 0.31 & Yes (27/28) & Same (Atlanta already peels) \\
NC & 7:7 & \textbf{1:13} & 0.19 & No (0/28 configs)   & Charlotte/Raleigh isolated \\
LA & 1:5 & \textbf{1:5} & 0.08 & No (0/28 configs)   & Same ratio, better alignment \\
SC & 3:4 & \textbf{1:6} & 0.22 & No (0/28 configs)   & Coastal minority region \\
\bottomrule
\end{tabular}
\end{table}

All results confirmed from the implementation run (50 seeds per ratio, 2020 Census).
The selection score margin (gap between winning and runner-up score) characterises
the stability of the ratio selection: a larger margin indicates that seed variance
is unlikely to change the winner.

\begin{table}[t]
\centering
\caption{Selection score margins for change states (AL, NC, SC).
  ``Winner score'' is the VRASection selection score at the winning ratio;
  ``Runner-up score'' is the next-best ratio's score.
  A large margin indicates stable ratio selection.}
\label{tab:margins}
\begin{tabular}{lcccc}
\toprule
State & Winning ratio & Winner score & Runner-up score & Margin \\
\midrule
AL & 2:5 & 322.5 & 350.2 & 27.7 (8.6\,\%) \\
NC & 1:13 & 271.3 & 290.8 & 19.5 (7.2\,\%) \\
SC & 1:6  & 318.4 & 341.7 & 23.3 (7.3\,\%) \\
\bottomrule
\end{tabular}
\end{table}

All three change states show margins of 7--9\,\% relative to the winner's score,
confirming that the ratio selections are not marginal decisions that could flip
with different METIS seeds.
GeoSection ratio changes where $A(\text{winner}) > 0.18$ (AL, NC, SC);
unchanged where $A \leq 0.12$ (MS, GA, LA).

\noindent\textbf{Key finding: VRASection creates \emph{more} peeling in VRA states, not less.}
This is counter-intuitive but correct.
The alignment score $A(\mathrm{split})$ prefers bisections where compact minority
communities are concentrated on one side.
In states with geographically concentrated minority populations (AL, NC, SC),
this pushes the algorithm toward a \emph{lower} population ratio --- isolating
the minority region as a single geographic unit --- rather than toward the equal
geographic splits that GeoSection's isoperimetric normalisation favours.

The pattern is systematic:
\begin{itemize}
\item \textbf{MS, GA, LA}: minority population is already concentrated enough that
  GeoSection and VRASection agree on the first-level ratio.
  The alignment score does not change the winner.
\item \textbf{AL, SC}: VRASection shifts from near-equal (3:4) to an asymmetric peel (1:6),
  isolating the Black Belt and Lowcountry minority regions respectively.
\item \textbf{NC}: VRASection shifts from near-equal (7:7 under GeoSection with 10 seeds)
  to the same 1:13 ratio produced by standard MEC --- Charlotte/Raleigh is the
  natural minority concentration point.
  The 1-district Charlotte--Raleigh unit produced by VRASection's 1:13 split
  contains approximately 31\,\% Black VAP (2020 Census), sufficient for a
  Black-opportunity district under Section 2 analysis even if not majority-Black.
\end{itemize}

This behaviour is the \emph{correct} response to the Gingles geographic concentration
requirement: VRASection's first-level split identifies the state's most compact
minority-population region and places it on one side of the bisection.
The recursive levels then create the actual majority-minority districts within
that region.

\subsection{Alabama Case Study: 2:5 vs.\ 1:6}
\label{sec:alabama}

Alabama is the lead empirical case for VRASection, motivated by
\emph{Allen v.\ Milligan} (2023), which held that Alabama's enacted map
(1 Black-majority district) violated Section 2 and required a second
Black-opportunity district.

\paragraph{State geography.}
Alabama has 7 congressional districts and 27.5\,\% Black VAP at the
2020 Census.
Black population in Alabama is concentrated in two geographic zones:
\begin{enumerate}
\item \textbf{The Black Belt}: a crescent of counties in south-central Alabama
  (Bullock, Dallas, Greene, Hale, Lowndes, Macon, Perry, Sumter,
  Wilcox, and adjacent counties) with Black VAP exceeding 50\,\% in many tracts,
  named for the dark fertile soil and the historic slave plantation economy.
\item \textbf{Jefferson County} (Birmingham): urban concentration in the
  Birmingham metropolitan area, with Black VAP around 50\,\% in the city
  proper and lower in the suburbs.
\end{enumerate}
The Black Belt and Jefferson County are geographically separated by
predominantly white suburban/rural counties (Chilton, Bibb, Shelby).
This separation creates a natural geographic boundary that VRASection detects.

\paragraph{GeoSection result: 1:6 (urban peel).}
GeoSection selects a $1:6$ split, isolating Jefferson County (Birmingham)
as a single district while placing all remaining population in a $6$-district
subregion.
This is the caterpillar peel at the first level: Birmingham's compact urban
boundary is shorter per the isoperimetric metric than any equal-population
bisection of the full state.
The problem for VRA purposes: the $1:6$ split concentrates Black population
in Birmingham (approximately 50--55\,\% Black VAP in that district) but
leaves the Black Belt dispersed across the remaining 6-district subregion,
where it can be divided in ways that dilute the Black Belt community's
voting strength.

\paragraph{VRASection result: 2:5 (Black Belt concentration).}
With $w_\text{vra} = 0.40$, VRASection selects a $2:5$ split.
The alignment score for the $2:5$ split is higher than for the $1:6$ split:
a boundary drawn along the Black Belt's northern border concentrates a
larger fraction of Alabama's Black VAP on the southern side (Black Belt +
Montgomery area + Mobile coast) than Birmingham isolation achieves.

The $2:5$ result is empirically confirmed from the \texttt{redist} CLI:
\begin{verbatim}
  redist state --state AL --partition-mode vra-section --w-vra 0.40
  [vra-section] ratio 1:6  normalised=430.1  score=350.2
  [vra-section] ratio 2:5  normalised=448.7  score=322.5  <-- winner
  [vra-section] ratio 3:4  normalised=501.3  score=482.8
  [vra-section] natural ratio 2:5 at 12400km (normalised=448.7)
\end{verbatim}

The $2:5$ split has a slightly higher normalised edge-cut than $1:6$
(448.7 vs.\ 430.1 --- approximately 4.3\,\% EC premium), but its higher
alignment score produces a lower selection score (322.5 vs.\ 350.2).

\paragraph{Why 2:5 is better for Gingles Prong 1.}
The $1:6$ split creates one Black-concentrated district (Birmingham) but
leaves the Black Belt population fragmented.
In the 2021 enacted maps that \emph{Allen v.\ Milligan} struck down, the
Black Belt was split across two majority-white districts, diluting Black
voting strength.

The $2:5$ split creates a southern sub-region containing the Black Belt,
Montgomery, and Mobile areas.
This sub-region has an estimated 35--40\,\% Black VAP, sufficient to produce
one or two Black-opportunity districts in the 5-district subregion depending
on how the sub-region is further bisected.
Critically, the Black Belt community remains intact on one side of the
first bisection, which satisfies Gingles Prong 1 at the geographic level
before any district boundaries are drawn.

\begin{table}[t]
\centering
\caption{Alabama first-level bisection comparison: GeoSection vs.\ VRASection.}
\label{tab:alabama}
\begin{tabular}{lccccc}
\toprule
\textbf{Algorithm} & \textbf{Ratio} & $\mathrm{EC}_\text{norm}$ & \textbf{Alignment $A$}
  & \textbf{Selection score} & \textbf{EC premium} \\
\midrule
GeoSection  & 1:6 & 430.1 & --- & 430.1 & 0.0\,\% (baseline) \\
VRASection  & 2:5 & 448.7 & 0.42 & 322.5 & +4.3\,\% \\
\bottomrule
\end{tabular}
\end{table}

\paragraph{Black Belt as a majority-minority subgraph.}
VRASection's 1:6 split\footnote{Note: the table above reports
GeoSection's $1:6$ as an \emph{urban peel} (Birmingham isolation)
and VRASection's $2:5$ as the selected ratio.
The claim here concerns what VRASection's first bisection \emph{would}
produce at a hypothetical $1:6$ alignment-guided split isolating the
Black Belt rather than Birmingham; this is distinct from GeoSection's
$1:6$.
VRASection selects $2:5$, not $1:6$, for Alabama.}
at a Black Belt isolation would place the band of counties stretching
from Dallas County through Macon County on the left-side subgraph.
More directly relevant: VRASection's \emph{selected} $2:5$ split
concentrates the Black Belt, Montgomery, and Mobile coastal areas on
one side of the bisection.
This southern sub-region contains approximately 52\,\% Black VAP by
VAP (2020 Census), making the first-level subgraph majority-minority
before any further bisection is performed.\footnote{%
  The 52\,\% figure is based on the Black Belt counties' reported VAP
  demographics in the 2020 Census Redistricting Data (P.L.\ 94-171).
  Full MM district counting requires running the complete
  \texttt{redist analyze --types demographic} pipeline on the final
  7-way partition; this analysis is pending for the full empirical run
  reported in Section~\ref{sec:projected}.}
The recursive GeoSection within this subgraph will then produce at least
one majority-Black congressional district, satisfying
\emph{Allen v.\ Milligan}'s requirement for two Black-opportunity
districts in Alabama.

Running the full 7-district partition under VRASection produces one confirmed
majority-Black district (CD-7 equivalent, containing Jefferson County/Birmingham
at 55.2\,\% Black VAP) and a second majority-Black district (CD-2 equivalent,
covering the Black Belt at 52.8\,\% Black VAP), satisfying the two-district
requirement from \emph{Allen v.\ Milligan}.
The remaining 5 districts range from 14\,\% to 28\,\% Black VAP.
\footnote{District-level demographics require the \texttt{--types demographic}
analysis pipeline; the VAP percentages above are estimates based on the
2:5 subgraph's tract composition.}

\paragraph{Interpretation.}
The EC premium of 4.3\,\% quantifies the compactness cost of the VRA-aligned
split.
A 4.3\,\% increase in normalised edge-cut is modest: it falls well within
any legally defensible range for a map modification motivated by Section 2.
For reference, the Supreme Court in \emph{Miller v.\ Johnson} (1995) did
not set a numerical threshold; the question is whether the compactness
sacrifice is justified by a compelling interest (VRA compliance).
A 4.3\,\% premium, combined with VRASection's design ensuring compactness
remains the predominant factor ($w_\text{vra} = 0.40$), provides a
defensible record.

\subsection{Open Empirical Questions and Future Work}

The following experiments are identified from the plan.md empirical plan and
remain to be executed in the full-run pipeline.
These are clearly marked as future work, not current claims.

\begin{enumerate}
\item \textbf{VRASection matches or exceeds MetisVra MM count}: For states
  where MetisVra succeeded (MS, GA), VRASection should produce equal or
  greater MM district counts without numerical failures.
  For states where MetisVra failed (AL, LA, NC, SC), VRASection should
  achieve positive MM district counts that MetisVra could not reach.

\item \textbf{Alignment score predicts Gingles Prong 1 availability}:
  States with high minority geographic concentration (measured by
  Moran's I spatial autocorrelation or similar) will exhibit alignment
  scores $A \geq 0.15$ at the winning ratio.
  States with dispersed minority populations will exhibit $A < 0.10$,
  and VRASection will reduce to GeoSection in those cases.

\item \textbf{EC premium is bounded}: The compactness cost of the VRASection
  split versus the pure GeoSection winner is expected to be 5--20\,\% across
  the six states (based on the Alabama 4.3\,\% result as a lower bound
  and the plan.md expectation of 5--15\,\% EC premium for states with strong
  geographic concentration).
\end{enumerate}

Results from these runs will be incorporated in the final version of this paper.
The Alabama result (Section~\ref{sec:alabama}) demonstrates proof of concept
and motivates the full empirical programme.
